\input{../tex-inputs/latex-leseansicht-vorspann}

\section[Emil Marriot an Arthur Schnitzler, 20. 6. 1902]{L04336 Emil Marriot an Arthur Schnitzler, 20. 6. 1902}
\nopagebreak\mylabel{L04336v}
\rehead{ }\normalsize\beginnumbering\briefempfaengerindex{Schnitzler, Arthur@\textsc{Schnitzler, Arthur}!zzzMarriot, Emil@\emph{von Emil Marriot}!1902-06-201@{20. 6. 1902}|(be}
\toendnotes[C]{\smallbreak\pagebreak[2]}
\correspDesc{Versand  durch Emil Marriot am 20. 6. 1902 in Brodzany
\newline{}Übermittlung  am 24. 6. 1902 in Brodzany
\newline{}Zustellung  am 25. 6. 1902 in Wien
\newline{}Erhalt  durch Arthur Schnitzler im Zeitraum [25. 6. 1902 –
                  26. 6. 1902] in Wien}\toendnotes[C]{\smallbreak}
\Standort{DLA, A:Schnitzler, HS.1985.1.4008.}
\physDesc{Bildpostkarte, 306 Zeichen
\newline{}Handschrift: schwarze Tinte, lateinische Kurrent
\newline{}Versand: 1) Stempel: »\nobreak{}\oindex{Brodzany@\textbf{Brodzany}|pwk}Brogyán, \textcolor{gray}{0}2 Jun 24\nobreak{}«.   2) Stempel: »\nobreak{}\oindex{IX., Alsergrund@\textbf{IX., Alsergrund}, \emph{Verwaltungsgebiet}|pwk}9/3 Wien 72, 25. 6. 02, 8.V, Bestellt\nobreak{}«. }\toendnotes[C]{\smallbreak}\pstart{}{\pb}Herrn\pend{}\pstart{}D\textsuperscript{r.} Arthur Schnitzler\pend{}\pstart{}Wien, IX\oindex{IX., Alsergrund@\textbf{IX., Alsergrund}, \emph{Verwaltungsgebiet}|pw}\pend{}\pstart{}Frankgasse 1\oindex{Wien@\textbf{Wien}!IX., Alsergrund@\textbf{IX., Alsergrund}!Frankgasse 1@\textbf{Frankgasse 1}, \emph{Wohngebäude}|pw}.\pend{}{\bigskip}
\pstart
           \noindent{}\centering{}{[}Ansicht des Haupthauses, des Gartens
                  und einer Kapelle von Schloss
                  Brodzany\oindex{Schloss Brodzany@\textbf{Schloss Brodzany}, \emph{Schloss}|pw}{]}\pend
           \vspace{1em}
\pstart
           \noindent{}{\pb}Aus\pend
           
\pstart
           \textcolor{gray}{\textbf{Schloss Brogyan}}\oindex{Schloss Brodzany@\textbf{Schloss Brodzany}, \emph{Schloss}|pw}\pend
           
\pstart
           \textcolor{gray}{\textbf{Barsmegye}}\oindex{Komitat Bars@\textbf{Komitat Bars}|pw}\pend
           
\pstart
           \textcolor{gray}{\textbf{Station Nagy-Bélicz}}\oindex{Bahnhof Veľké Bielice@\textbf{Bahnhof Veľké Bielice}, \emph{Bahnhofsgebäude}|pw} in Ungarn\oindex{Ungarn@\textbf{Ungarn}|pw} sendet einen kollegialen
               Gruss\pend
           
\pstart
           Ihre ganz ergebene{\\[\baselineskip]}\spacefill\mbox{Marriot.}\pend
           \leftskip=0em{}
\pstart
           20. Juni 1902\pend
           
\pstart
           \noindent{}Sind sie überhaupt in Wien\oindex{Wien@\textbf{Wien}, \emph{Verwaltungsgebiet}|pw}? Hier ist viel
                  »Stimmung« – und doch \label{K_L04336-1v}\edtext{fällt mir
                  nichts ein}{\lemma{\textnormal{\emph{fällt mir
                  nichts ein}}}\Cendnote{\textnormal{Marriot\pwindex{Marriot, Emil 20.\,11.\,1855 Wien – 5.\,5.\,1938 ebd.@\textsc{Marriot, Emil} (20.\,11.\,1855 Wien – 5.\,5.\,1938 ebd.), \emph{Schriftstellerin}|pwk} verarbeitete den Aufenthalt auf dem Schloss in
                     ihrem Roman \emph{Anständige Frauen}\pwindex{Marriot, Emil 20.\,11.\,1855 Wien – 5.\,5.\,1938 ebd.@\textsc{Marriot, Emil} (20.\,11.\,1855 Wien – 5.\,5.\,1938 ebd.), \emph{Schriftstellerin}!Anständige Frauen. Roman@\strich\emph{Anständige Frauen. Roman}|pwk}
                        (1906).}}}\label{K_L04336-1}. Aber schon gar nichts!\pend
           \selectlanguage{ngerman}\endnumbering\briefempfaengerindex{Schnitzler, Arthur@\textsc{Schnitzler, Arthur}!zzzMarriot, Emil@\emph{von Emil Marriot}!1902-06-201@{20. 6. 1902}|)be}\mylabel{L04336h}
\begin{anhang}
\end{anhang}\newcommand{\dateiname}{L04336}\newcommand{\titel}{Emil Marriot an Arthur Schnitzler, 20. 6. 1902}\newcommand{\editorInnen}{Herausgegeben von Jahnke, SelmaMüller, Martin Anton}\input{../tex-inputs/latex-leseansicht-abspann}
